\chapter{绪 论}

“什么是黎曼几何学？”每一位初学者在打开本书时都会提出这样的问题. 对这个问题的回答既是简单、容易的, 又是复杂、困难的. 让我们从 Gauss 的 “绝妙定理” (Theorema Egregium) 谈起.

为了刻画三维欧氏空间中正则参数曲面的形状，通常要引进曲面的第一基本形式和第二基本形式的概念。第一基本形式是曲面上的切向量  $\mathrm{d}\vec{r}$  的长度平方，即

$$
I = \mathrm {d} \vec {r} \cdot \mathrm {d} \vec {r}.
$$

第二基本形式是

$$
I I = \mathrm {d} ^ {2} \vec {r} \cdot \vec {n}
$$
(其中  $\vec{n}$  为曲面的单位法向量), 在本质上它是曲面上任意一点的邻近点到该点切平面的有向距离. 特别是, 两个基本形式之比

$$
\frac {\varPi}{I} = \frac {\mathrm {d} ^ {2} \vec {r} \cdot \vec {n}}{\mathrm {d} \vec {r} \cdot \mathrm {d} \vec {r}} = \frac {\mathrm {d} ^ {2} \vec {r}}{\mathrm {d} s ^ {2}} \cdot \vec {n}
$$
是曲面上通过该点, 以  $\mathrm{d}\vec{r}$  为方向的曲线的曲率向量  $\kappa \vec{\beta} = \frac{\mathrm{d}^2\vec{r}}{\mathrm{d}s^2}$  在曲面该点处的单位法向量  $\vec{n}$  上的投影 (这里的  $\kappa$  和  $\vec{\beta}$  分别是曲线的曲率和主法向量), 它是仅依赖曲面在该点的切方向的函数. 如果考虑曲面上由该点的切方向  $\mathrm{d}\vec{r}$  和单位法向量  $\vec{n}$  所张成的平面, 并且用该平面在曲面上截出一条曲线, 则这条曲线以  $\mathrm{d}\vec{r}$  为切向量, 同时它在该点的曲率向量与  $\vec{n}$  是共线的. 所以该曲线在该点的曲率正好是  $\frac{II}{I}$ , 我们把它称为曲面在该点沿切方向  $\mathrm{d}\vec{r}$  的法曲率, 记为  $\kappa_{n}$ . 在曲面上任意固定一点, 则  $\kappa_{n}$  是在该点的切方向的函数, 它反映了曲面在该点沿该切方向的弯曲方向和弯曲程度. 一般说来, 曲面在每一点有两个彼此垂直的切方向, 使得法曲率  $\kappa_{n}$  在这两个方向分别达到它的最大值和最小值.

这两个切方向称为曲面在该点的主方向, 相应的两个法曲率称为曲面在该点的主曲率, 记为  $\kappa_{1}, \kappa_{2}$ . 所谓的 Gauss 曲率  $K$  指的就是这两个主曲率的乘积, 即  $K = \kappa_{1} \kappa_{2}$ . 自然, 它是借助于曲面的第一基本形式和第二基本形式计算而得的.

Gauss 经过复杂的计算, 获得了一个惊人的发现 (1827 年): Gauss 曲率  $K$  只依赖于曲面的第一基本形式, 而与曲面的第二基本形式无关. 这就是 Gauss 的绝妙定理.

Gauss 的绝妙定理的意义在哪里？如果我们把参数曲面的定义域记为  $D$ ，它是  $\mathbb{R}^2$  中的一个开子集，其中的点的坐标记为  $(u^1, u^2)$ ，那么曲面的第一基本形式  $I$  是在区域  $D$  上坐标  $u^1, u^2$  的 2 次微分式：

$$
I = \sum_ {i, j = 1} ^ {2} g _ {i j} (u ^ {1}, u ^ {2}) d u ^ {i} d u ^ {j},
$$
其中  $g_{ij} = g_{j\prime}$  ，并且在每一点  $(u^{1},u^{2})\in D,(g_{ij})$  是正定的  $2\times 2$  矩阵.在这样一种结构下，能够做些什么事呢？按照曲面论，我们能够计算区域  $D$  上任意一条分段光滑曲线的长度，能够计算区域  $D$  内一个有界子区域的面积等等；而且这些量与定义区域  $D$  内所取的坐标系无关.Gauss的定理则进一步断言：利用  $(g_{ij})$  ，我们还可以计算曲面在  $D$  内每一点处的Gauss曲率，而不管曲面的具体形状如何.这是一个非常了不起的结果，它开创了曲面的内蕴微分几何.Gauss曲率的意义是什么？它所反映的不只是我们所观察到的曲面的“外在”形状，而且是衡量定义在区域  $D$  上的第一基本形式  $I$  与标准的欧氏度量偏离程度的量度.以Gauss曲率  $K$  为常数  $c$  的第一基本形式为例(参看参考文献[2],第184页):

当  $c = 0$  时,  $I = (\mathrm{d}u^{1})^{2} + (\mathrm{d}u^{2})^{2}$ ;

当  $c > 0$  时,  $I = (\mathrm{d}u^{1})^{2} + \cos^{2}(\sqrt{c} u^{1})(\mathrm{d}u^{2})^{2}$ ;

当  $c < 0$  时,  $I = (\mathrm{d}u^{1})^{2} + \cosh^{2}(\sqrt{-c} u^{1})(\mathrm{d}u^{2})^{2}$ .

通过坐标变换, 可以用等温参数把上述三种情况统一起来:

$$
I = \frac {(\mathrm {d} u ^ {1}) ^ {2} + (\mathrm {d} u ^ {2}) ^ {2}}{[ 1 + \frac {c}{4} ((u ^ {1}) ^ {2} + (u ^ {2}) ^ {2}) ] ^ {2}}.
$$
更一般地, 可以证明: 在  $D$  上的第一基本形式  $I$  是欧氏度量, 当且仅当它的 Gauss 曲率  $K$  恒为零. 由此可见, Gauss 曲率衡量了第一基本形式相对于欧氏度量的偏离程度, 即曲面 (或二维空间) 的 “内在” 弯曲程度. 我们可以研究各种不同的二维 “弯曲” 空间, 而经典的非欧几何恰恰是在常 “弯曲” 空间中的几何学,

B. Riemann 在 1854 年给出的著名就职演说《关于几何学的基本假设》(参看参考文献 [30], Vol.2, 第 135 页) 中把 Gauss 的曲面内蕴微分几何推广到任意维数的情形。首先，他提出了  $n$  维流形的概念。尽管在当时，尚没有具体、确切地表述度量空间和拓扑空间等概念的方式，但在他的头脑里已把  $n$  维流形设想为在局部上与  $n$  维欧氏空间相仿的对象，其中每一个点都可以用  $n$  个有序实数的组  $x = (x^{1}, \dots, x^{n})$  来描写。在欧氏空间中，为了求曲线的长度，需要先指定直线段的长度，然后把曲线的长度定义为它的内接折线长度的上确界。在这里，“直线段”是一类特殊的曲线，需要预先把它从一般的曲线中区分出来。现在，Riemann 提出一种与之不同的统一的模式，不需要预先区分直线和曲线，而是先定义切向量的长度，然后把曲线的长度定义为切向量的长度沿曲线的积分。因而，这种做法适用于任意的（光滑）流形，不必要求该流形有如同欧氏空间的平直结构。Riemann 进一步提出：切向量  $\mathrm{d}x$  的长度  $\mathrm{d}s$ （称为线元）可以是  $\mathrm{d}x$  的分量的任意的一次齐次函数，要求该函数的值在  $\mathrm{d}x$  的分量全部反号时不改变；并且该函数的系数与  $x$  有关（后来，Finsler 首次对此作了系统的研究，现在称这种度量为 Finsler 度量）。特别地， $\mathrm{d}s$  可以是  $\mathrm{d}x^{i}$  的、处处为正的二次齐次函数的算术平方根，而其中的系数是变量  $x$  的连续函数。后者就是现在的 Riemann 度量。由此可见，简单地说，黎曼几何恰恰是 Gauss 的曲面内蕴微分几何在高维的推广，曲面的内蕴微分几何是二维黎曼流形的几何学。

Gauss 对 Riemann 的演讲评价很高, 他在步出演讲厅时以罕有的激动心情对 W. Weber 谈起 “Riemann 所表述的观点的深度” (参看参考文献 [30], Vol.2, 第 134 页). 事实的确是如此, 为了理解并在技术细节上完善 Riemann 的思想差不多花了大半个世纪. 首先是 Christoffel,

然后是 Ricci 和 Levi-Civita, 他们创造了一整套张量分析的方法, 引进了所谓的绝对微分学, 给出了 Riemann 所提出的曲率的表达式. 尤其是 Levi-Civita 提出了曲面上的切向量沿曲线平行移动的概念. 这样,对于黎曼几何在几何直观上的理解便提高到一个崭新的水平, 大大地推动了黎曼几何的发展.

1916年,A.Einstein在他所发表的广义相对论中成功地运用了黎曼几何学, 把质量分布表述为黎曼度量, 把引力现象解释为黎曼空间的曲率性质. 于是, 黎曼几何开始受到普遍的重视, 研究“弯曲”空间的必要性再次得到肯定. 在其后不久出版的教科书:

L. P. Eisenhart, Riemannian Geometry, Princeton University Press, New Jersey, 1926

在传播黎曼几何知识方面起到了重大的作用. 与此同时, E. Cartan 出版了他的重要著作:

E. Cartan, Lecons sur la Géométrie des Espaces de Riemann, Gauthier-Villars, Paris, 1928.

流形论是黎曼几何的基础。对流形（拓扑流形）的概念第一次做出准确描述的是 D. Hilbert（《几何基础》，1902 年）。后来，H. Weyl 在他的名著《黎曼面的概念》(1913 年) 中对于微分流形给出了清晰的数学描述。在 20 世纪 30 年代，H. Whitney 开始对微分流形的拓扑进行了认真的研究。在做了这些准备之后，H. Hopf 在 1932 年提出了研究大范围黎曼几何的问题，即截面曲率  $K$  的符号与（紧致）黎曼流形  $(M, g)$  的拓扑相互制约的问题（参看 M. Berger, Riemannian Manifolds: From Curvature to Topology, In: Chern — A Great Geometer of the Twentieth Century, 第 184 页, International Press, Hong Kong, 1992）。在 1942 年，陈省身用 Cartan 的外微分和活动标架方法，成功地用内蕴的方法证明了偶数维紧致黎曼流形上的 Gauss-Bonnet 定理，这是大范围黎曼几何发展过程中的里程碑，开创了大范围黎曼几何的新纪元。陈省身在他的论文：A simple intrinsic proof of the Gauss-Bonnet formula for closed Riemannian manifolds (Ann. of Math., 45(1944), 747-752) 中以深邃的思想揭示了联络、曲率形式、切丛和球丛的重要性以及它们之间的相

互联系, 编织了黎曼流形的局部不变量 (曲率) 和整体不变量 (Euler 示性数) 之间的辉煌图景。陈省身的这项研究成果以及随后发表的关于陈示性式的一系列论文, 展示了后来的微分几何发展方向, 影响极为深远, 很快使得微分几何的思想和方法成为拓扑学、代数几何学、大范围分析、数学物理等许多数学分支的有机组成部分, 并使微分几何成为数学研究的一个中心课题。

在这里, 我们不能不提到陈省身在芝加哥大学讲授微分几何的油印讲义 Differentiable Manifolds (1953 年), 它以清晰的现代语言讲述微分流形、联络、黎曼流形等基本理论. 后来出版的微分流形与黎曼几何教科书都受到这本讲义的深刻影响. 当代许多重要的微分几何学家都是在陈省身的研究工作及上述讲义的养育、鼓舞下成长起来的. 陈省身的著作成为当代微分几何学家的必读经典.

从黎曼几何的发展过程, 以及黎曼几何在各个数学分支中的应用来看, 联络的概念处于中心的位置. 在流形上给定微分结构之后, 微分流形上的切向量、切空间和光滑切向量场都是有意义的, 因为它们所涉及的是微分流形上的光滑函数, 以及光滑函数的导数. 若要在微分流形上进一步运用微分手段, 必须能够对光滑切向量场求微分. 然而, 微分流形的光滑结构本身并未提供对光滑切向量场求导的手段. 对光滑切向量场求导本身是加在光滑流形上的一种新的结构, 这种结构就是所谓的联络. 有幸的是, 在黎曼流形上存在一种自然的特殊联络 (称为黎曼联络或Levi-Civita联络), 它是被黎曼结构 (即黎曼度量) 唯一确定的. 由此可见, 联络是比黎曼度量更为基本的一种结构. 在20世纪的五六十年代, 围绕联络的概念进行了紧张的研究, 其主要结晶是建立并完善了向量丛上的联络 (Koszul) 和主丛上的联络 (Ehresmann) 等概念. 陈示性式是通过复向量丛上的联络构造出来的, 然而它是与向量丛上联络的选取无关的不变量 (Chern-Weil 定理), 因此它反映了复向量丛的结构特性, 特别是量度了复向量丛偏离平凡丛 (空间的直积)的程度.

从20世纪五六十年代以来，大范围黎曼几何本身及其在各个数学分支中的应用已经发展到相当可观的程度。例如，度量的曲率性质与

流形本身的拓扑的相互制约关系, de Rham-Hodge 理论, Yang-Mills 理论, 黎曼流形的谱, 调和映射, 黎曼流形上的函数论, 黎曼流形中的极小子流形, 热流方法和问题, 非线性理论中的几何问题, Gromov 的黎曼度量的收敛问题, 等等, 都是黎曼几何中的重要课题. 由于黎曼几何已经有如此众多的分支, 所以, 我们说, “什么是黎曼几何学?” 是一个难以确切回答的问题. 在本书最后所列的参考书目, 有些是新近出版的黎曼几何教科书, 有些是黎曼几何专著, 往往侧重于或强调黎曼几何的某些课题. 若想更多地了解上述问题的答案, 读者可以选读这些书

“在本课程中能够学到些什么？”本课程是黎曼几何学的入门课程。其先修课程是“微分流形”。本课程面向基础数学专业、应用数学专业（以及理论物理专业和近代力学专业）的全体硕士研究生和博士研究生，主要目标是向他们介绍黎曼几何的基本概念和基础理论，大体上可以分为三个部分。鉴于联络的重要性，我们首先在联络的引进、联络的一般概念和黎曼联络的特性、切向量的平行移动等方面倾注了很大的力量，务使读者对此有一个清晰的了解。然后在黎曼流形上对测地线、弧长第一变分公式和曲率展开讨论。前四章是黎曼几何的基础，它们构成本书（上册）的第一部分。在第二部分，我们要导出弧长的第二变分公式，在此基础上，讨论测地线的最短性（即短程性），Jacobi场和共轭点理论，以及它们在大范围黎曼几何中的应用。这部分内容的教学目的是帮助读者初步建立起大范围黎曼几何的观念，并掌握研究大范围黎曼几何的变分方法。第三部分旨在建立黎曼子流形的框架理论，并导出子流形体积的第一、第二变分公式。

除了前七章作为标准的“黎曼几何引论”课程的基本内容以外，我们还辟专章介绍Kähler流形、黎曼对称空间等特殊黎曼流形，以及主丛上的联络。它们构成本书的下册。这些内容本身是微分几何的重要研究课题，而且是黎曼几何以及微分几何在各个数学分支中的应用的重要的基础。当然，在这里我们强调的还是这些内容的基础理论，而不是它们的最新发展。但是，这些基础理论对于从事有关研究工作的读者来讲是必备的知识。